\documentclass[10pt,english]{beamer}


\usetheme[math]{UBC}

\setbeamertemplate{caption}[numbered]

% Tables
\usepackage{booktabs}

% Figures & Media 
\usepackage{graphicx}
\usepackage{tikz}

% Metadata
\usepackage{pdfpc}
\usepackage{pdfpc-movie}

% Bibliography
\usepackage[
    backend=biber,
    style=numeric,
    sorting=none,
    maxbibnames=3,
    giveninits=true,
    doi=true,
    url=false,
    isbn=false,
]{biblatex}

\graphicspath{
    {figure/vector/}
        {figure/raster/}
}

% Math 
\usepackage{amsmath}
\usepackage{bm}
\usepackage{derivative}

\newcommand{\Reynolds}{\mathrm{Re}}
\newcommand{\Rep}{\mathrm{Re}_p}
\newcommand{\Fr}{\mathrm{Fr}}
\newcommand{\St}{\mathrm{St}}
\newcommand{\KB}{K_B}
\newcommand{\dd}{\mathrm{d}}
\newcommand{\transp}{\mathsf{T}}
\newcommand{\norm}[1]{\left\lVert#1\right\rVert}
\let\vec\undefined
\newcommand{\vec}[1]{\bm{#1}}
\let\epsilon\undefined
\newcommand{\epsilon}{\varepsilon}

\newcommand{\EBpos}{\vec X}
\newcommand{\EBaug}{r}
\newcommand{\EBlagmult}{\vec T}
\newcommand{\EBflex}{EI}
\newcommand{\EBlindens}{\rho_b}

\newcommand{\IBMforcedens}{\vec\Lambda}
\newcommand{\IBMforcelag}{\vec F}
\newcommand{\IBMforceeul}{\vec f}

\addbibresource{references.bib}
\renewcommand*{\bibfont}{\small}
\renewcommand*{\mkbibacro}[1]{#1}
\setbeamertemplate{bibliography item}{\insertbiblabel}

\title[Flexible Fibres on Adaptive Octree Grids]{Simulation of Flexible Filaments in Inertial Flow on Adaptive Octree Grids}
\subtitle{}
\author{Conor Olive\inst{1} \and Anthony Wachs\inst{1,2} \and Mona Rahmani\inst{2}}
\institute{\inst{1} University of British Columbia, Mathematics \and \inst{2} University of British Columbia, Chemical and Biological Engineering}
\date{\today}

\begin{document}

\maketitle

\begin{frame}{Motivation}
    \begin{itemize}
        \item Flexible fibres occur in various natural and anthropogenic
              contexts
        \item Samples of wastewater/effluent (especially near large North
              American urban centres) indicate that plastic fibres make up the
              largest share of microplastic type
              \cite{sutton_microplastic_2016}
        \item Resolved numerical studies of many-fibre aggregate dynamics in
              inertial regimes are sparse \cite{du_roure_dynamics_2019}
    \end{itemize}
    \vfill
    \begin{alertblock}{Problem Statement}
        How can we build a reasonably resolved numerical model to study
        \textbf{many-fibre}, potentially very \textbf{dilute} suspentions in
        various \textbf{inertial} flow regimes?
    \end{alertblock}
\end{frame}

\begin{frame}
    \begin{figure}
        \centering
        \includegraphics[width=0.9\textheight]{microplastic-types.png}
        \caption{Different categories of microplastics: (A) fragment, (B) film, (C) foam, (D) fibre, (E) line, and (F) pellet. Reproduced from \textcite{browne_accumulation_2011}.}
        \label{fig:microplastic-types}
    \end{figure}
\end{frame}

\begin{frame}{Continuous Model}
    \small
    \begin{columns}[T,onlytextwidth]
        % \begin{itemize} \item We choose to model the microplastic fibres as
        % large-angle Euler--Bernoulli (EB) beams and fluid with Navier--Stokes
        % equations \item Spectial treatement of the discretization and
        % fluid-structure coupling is required to make large problems
        % computationally tractable \end{itemize} \end{column}
        \begin{column}{0.5\textwidth}
            \begin{block}{Incompressible Navier-Stokes}
                \begin{minipage}[t][2.25cm][t]{\linewidth}
                    \begin{align*}
                        \rho_f \left(\pdv{\vec u}{t} + \vec u \cdot \nabla \vec u\right) & =
                        - \nabla p + \mu_f \nabla^2 \vec u
                        + \rho_f \vec g + \IBMforceeul                                         \\
                        \nabla \cdot \vec u                                              & = 0
                    \end{align*}
                \end{minipage}
            \end{block}
        \end{column}
        \begin{column}{0.475\textwidth}
            \begin{block}{Large Angle Euler-Bernoulli Beam}
                \begin{minipage}[t][2.25cm][t]{\linewidth}
                    \begin{align*}
                        \EBlindens \pdv[2]{\EBpos}{t}
                        + \EBflex \pdv[4]{\EBpos}{s}
                        - \pdv{}{s} \left(\EBlagmult \pdv{\EBpos}{s}\right) & =
                        \IBMforcelag + \EBlindens\vec g                           \\
                        \norm{\pdv{\EBpos}{s}}^2                            & = 1
                    \end{align*}
                \end{minipage}
            \end{block}
        \end{column}
    \end{columns}
    \vspace{0.6em}
    \begin{block}{Immersed Boundary}
        \begin{equation*}
            \begin{aligned}
                \pdv{\EBpos}{t}(s,t)
                 & = \vec u(\EBpos(s,t),t),
                \qquad
                \IBMforcelag(s,t)=-\IBMforcedens(s,t), \\
                \IBMforceeul(\vec x,t)
                 & = \int_0^L \IBMforcedens(s,t)\,
                \delta(\vec x-\EBpos(s,t))\,\odif{s} .
            \end{aligned}
        \end{equation*}
    \end{block}
\end{frame}

\begin{frame}{Discretization}
    \begin{itemize}
        \item An explicit direct-forcing immersed-boundary method is chosen to
              model the fluid-structure interaction in a partitioned manner,
              where the force density is obtained from a kinematic no-slip
              constraint
        \item The open source Basilisk \cite{popinet_quadtree-adaptive_2015}
              solver is chosen to advance the fluid equations, which provides
              tree-structured adaptive mesh refinement and an off-the-shelf
              fractional-step N-S solver
        \item The fibre is discretized using an augmented Lagrangian method
              first described by \textcite{glowinski_large_nodate}, needed to
              treat the inextensibility constraint
    \end{itemize}
\end{frame}

\begin{frame}
    \begin{figure}
        \includegraphics[width=0.9\textwidth]{octree.png}
    \end{figure}
\end{frame}

\begin{frame}{Augmented-Lagrangian Beam Solve}
    \begin{itemize}
        \item Alternating minimization with respect to \(\vec p\) and
              \(\EBpos\), followed by a multiplier update, yields an
              alternating direction method enforcing inextensibility
        \item The \(\vec p\)-subproblem remains local in \(s\) while the
              \(\EBpos\) subproblem is a linear fourth-order beam solve using
              the current multiplier and tangent field fixed
        \item The linear system resulting from finite-element discretization is
              sparse, symmetric, and positive definite, and remains constant
              for fixed augmentation parameters and appropriate B.C.s
        \item A single factorization can be reused for all components and all
              alternating-direction iterations
        \item A average-acceleration Newmark method is used to discretize the
              time derivative, where the static subproblem is solved at each
              time step
    \end{itemize}
\end{frame}

\begin{frame}
    \pdfpcnote{
        Here p is an auxiliary tangent field.
        It approximates the first spatial derivative dX/ds.
        The constraint p = dX/ds is enforced by the multiplier update.
    }
    Given \(\EBpos_h^0\), \(\EBlagmult_h^0\), penalty \(\EBaug>0\), repeat:
    \begin{enumerate}
        \item Project the tangent field:
              \[
                  \vec p_h^{k+1}
                  =
                  \frac{
                      \EBaug\,\partial_s \EBpos_h^k - \EBlagmult_h^k
                  }{
                      \norm{\EBaug\,\partial_s \EBpos_h^k - \EBlagmult_h^k}
                  }.
              \]

        \item Solve the linear beam problem for \(\EBpos_h^{k+1}\):
              \[
                  (EI\,K_b+\EBaug K_s)\EBpos_h^{k+1}
                  =
                  \vec F(\vec p_h^{k+1},\EBlagmult_h^k).
              \]

        \item Update the multiplier:
              \[
                  \EBlagmult_h^{k+1}
                  =
                  \EBlagmult_h^k
                  +
                  \alpha
                  \left(
                  \vec p_h^{k+1}
                  -
                  \partial_s\EBpos_h^{k+1}
                  \right).
              \]

        \item Stop when
              \[
                  \norm{\vec p_h^{k+1}-\partial_s\EBpos_h^{k+1}} < \epsilon
              \]
    \end{enumerate}
\end{frame}

\begin{frame}{Beam Validation}
    \begin{columns}[T,onlytextwidth]
        \begin{column}{0.49\textwidth}
            \centering
            \input{figure/vector/addm-static-convergence.tex}

            {\small Static Bisshopp--Drucker cantilever
                \parencite{bisshopp_large_1945}}
        \end{column}
        \begin{column}{0.49\textwidth}
            \centering
            \input{figure/vector/addm-nonlinear-convergence.tex}

            {\small Nonlinear inextensible manufactured solution }
        \end{column}
    \end{columns}
\end{frame}

\begin{frame}{Direct-Forcing Immersed Boundary Method}
    % Given predicted fluid velocity \(\vec u^*\) and fibre state, for each time
    % step:
    \begin{enumerate}
        \item Interpolate fluid velocity to Lagrangian fibre markers:
              \[
                  \vec U_a^* = J\vec u^* .
              \]

        \item Compute the target marker velocity from the beam update:
              \[
                  \vec U_a^{\,\mathrm{target}}
                  =
                  \dot{\vec X}_a^{\,n+1}.
              \]

        \item Choose the Lagrangian force density \(\IBMforcedens \in
              \mathbb{R}^d\) from the slip residual:
              \[
                  \IBMforcedens
                  =
                  \frac{\rho_f}{\Delta t}
                  \left(
                  \vec U_a^{\,\mathrm{target}}
                  -
                  \vec U_a^*
                  \right).
              \]

        \item Spread the marker force back to the Eulerian grid, and use this
              to correct fluid momentum
              \[
                  \IBMforceeul = S\IBMforcedens, \qquad
                  \vec u^{**}
                  =
                  \vec u^*
                  + \frac{\Delta t}{\rho_f} \IBMforceeul .
              \]
    \end{enumerate}
\end{frame}

\begin{frame}{Direct-Forcing Immersed Boundary Method Validation}
    \begin{columns}[T,onlytextwidth]
        \begin{column}{0.5\textwidth}
            \input{figure/vector/static-sphere-Re-sweep.tex}
        \end{column}
        \begin{column}{0.5\textwidth}
            \begin{block}{Schiller-Naumann}
                \[
                    C_D(\Rep)
                    =
                    \frac{24}{\Rep}
                    \left(1+0.15\Rep^{0.687}\right)
                \]
            \end{block}

            \begin{itemize}
                \item Uniform flow past a fixed sphere.
                \item IBM force is integrated as drag.
                \item Tested over \(\Rep=10\)--\(250\).
            \end{itemize}
        \end{column}
    \end{columns}
\end{frame}

\begin{frame}{Fully Coupled Validation}
    \begin{columns}
        \centering
        \begin{column}{0.85\textwidth}
            \centering
            \href{run:figure/raster/single-fibre-Re-200-KB-10e-3_surface_edges.mp4?autostart}{
                \includegraphics[width=\textwidth]{single-fibre-Re-200-KB-10e-3_surface_edges_0001.png}
            }
        \end{column}
    \end{columns}
\end{frame}

\begin{frame}{}
    \begin{center}
        \centering
        \begin{tabular}{@{}cccc@{}}
            \includegraphics[height=0.85\textheight]{single-fibre-Re-200-KB-10e-3_early_t9p2.png}  &
            \includegraphics[height=0.85\textheight]{single-fibre-Re-200-KB-10e-3_early_t10.png}   &
            \includegraphics[height=0.85\textheight]{single-fibre-Re-200-KB-10e-3_early_t10p8.png} &
            \includegraphics[height=0.85\textheight]{single-fibre-Re-200-KB-10e-3_early_t11p6.png}   \\
            {\scriptsize \(t=9.2\)}                                                                &
            {\scriptsize \(t=10.0\)}                                                               &
            {\scriptsize \(t=10.8\)}                                                               &
            {\scriptsize \(t=11.6\)}
        \end{tabular}
        \par\vspace{0.4em}
        {\scriptsize \(\Reynolds=200\), \(\KB=10^{-3}\). Black curves show
        vorticity contours; red markers show Lagrangian fibre points.}
    \end{center}
\end{frame}

\begin{frame}
    \begin{columns}
        \begin{column}{5in}
            \small
            \centering
            \input{figure/vector/huang-tip-Re-200-KB-10e-3.tex}
            \setlength{\tabcolsep}{3pt}
            \begin{tabular}[t]{@{}lc@{}}
                \toprule
                Parameter     & Value       \\
                \midrule
                \(\Reynolds\) & \(200\)     \\
                \(q\)         & \(1.5\)     \\
                \(\KB\)       & \(10^{-3}\) \\
                \(\Fr\)       & \(0.5\)     \\
                \bottomrule
            \end{tabular}
            \hspace{0.25in}
            \begin{tabular}[t]{@{}lcc@{}}
                \toprule
                                                          & Amplitude             & \(\St\)           \\
                \midrule
                Huang et al. \cite{huang_simulation_2007} & \(\pm 0.35\)          & \(0.30\)          \\
                Wang \cite{wang_strongly_2015}            & \(\pm 0.35\)          & \(0.31\)          \\
                Lee \cite{lee_discrete-forcing_2015}      & \(\pm 0.38\)          & \(0.31\)          \\
                Goza \cite{goza_strongly-coupled_2016}    & \(\pm 0.38\)          & \(0.32\)          \\
                \textbf{Present}                          & \(\mathbf{\pm 0.36}\) & \(\mathbf{0.31}\) \\
                \bottomrule
            \end{tabular}
        \end{column}
    \end{columns}
\end{frame}

\begin{frame}{Conclusion}
    \begin{columns}[T,onlytextwidth]
        \begin{column}{0.48\textwidth}
            \begin{block}{Current Status}
                \begin{minipage}[t][0.38\textheight][t]{\linewidth}
                    \begin{itemize}
                        \item Individual partitioned models show good
                              validation
                        \item Fully coupled fibre motion agrees with published
                              results
                    \end{itemize}
                \end{minipage}
            \end{block}
        \end{column}
        \begin{column}{0.48\textwidth}
            \begin{block}{Future Work}
                \begin{minipage}[t][0.38\textheight][t]{\linewidth}
                    \begin{itemize}
                        \item Complete additional validation cases in the short
                              term
                        \item Extend the model with additional physical
                              mechanisms
                        \item Add contact model(s) to extend model for dense
                              fibre packings
                    \end{itemize}
                \end{minipage}
            \end{block}
        \end{column}
    \end{columns}

    \vfill
    Contact me: \href{mailto:colive@math.ubc.ca}{colive@math.ubc.ca}
\end{frame}

\begin{frame}[plain,noframenumbering]
    \centering
    \vfill
    {\usebeamerfont{section title}\usebeamercolor[fg]{section title}
        Questions?
    }
    \vfill
\end{frame}

\begin{frame}[allowframebreaks]{References}
    \printbibliography[heading=none]
\end{frame}

\end{document}
